% Multi-Field Perturbation Theory: Architecture Design
% Source: docs/multi_field_perturbation.md (migrated to LaTeX)
% Uses macros from preamble.tex

\section{Multi-Field Perturbation Theory: Architecture Design}
\label{sec:multi-field-perturbation}

\subsection{Motivation}
\label{sec:mfp-motivation}

The current \texttt{[linearization]} pipeline handles a single perturbation field --- the metric:
\begin{equation}
  \gmunu = \etamunu + \varepsilon\, h_{\mu\nu}
  \label{eq:mfp-metric-pert}
\end{equation}
via \texttt{SetupMetricPerturbation} + \texttt{Perturbation[L, 2]} + \texttt{VarD[h]}.

Many physically interesting theories require \textbf{multiple perturbation fields} expanded simultaneously around their backgrounds. The Gertsenshtein effect (Einstein--Maxwell) is the immediate use case, but the architecture is general-purpose.

\subsection{Use Cases}
\label{sec:mfp-use-cases}

\begin{table}[htbp]
\centering
\caption{Example theories requiring multi-field perturbation.}
\label{tab:mfp-use-cases}
\begin{tabular}{@{}lll@{}}
\toprule
Theory & Perturbation Fields & Background \\
\midrule
\textbf{Einstein--Maxwell} (Gertsenshtein) & metric $h_{\mu\nu}$ + EM vector $a_\mu$ & $\etamunu + \bg{A}_\mu$ (uniform $\Bzero$) \\
\textbf{Einstein--Proca} & metric $h_{\mu\nu}$ + massive vector $a_\mu$ & $\etamunu$ \\
\textbf{Einstein--Yang--Mills} & metric $h_{\mu\nu}$ + gauge field $a_\mu^I$ & $\etamunu + \bg{A}_\mu^I$ \\
\textbf{Einstein--Cartan} & metric $h_{\mu\nu}$ + torsion $T^\lambda{}_{\mu\nu}$ & $\etamunu + \bg{T}$ \\
\textbf{Scalar-tensor gravity} & metric $h_{\mu\nu}$ + scalar $\varphi$ & $\etamunu + \bg{\varphi}$ (VEV) \\
\textbf{Matter-only} (fixed geometry) & vector $a_\mu$ (or scalar, tensor) & $\etamunu$ (fixed) + $\bg{A}_\mu$ \\
\bottomrule
\end{tabular}
\end{table}

\subsection{TOML Configuration}
\label{sec:mfp-toml}

\subsubsection{Full Multi-Field Perturbation (e.g., Einstein--Maxwell)}
\label{sec:mfp-toml-full}

\begin{lstlisting}[style=toml]
[linearization]
perturbation_field = "h"   # Metric perturbation (existing, unchanged)

[[linearization.matter_perturbations]]
field = "A"                # Must match a [[fields]] entry
perturbation_name = "a"    # Name for the 1st-order perturbation tensor
background = "Abar"        # Resolved via matching [[background_fields]] entry
\end{lstlisting}

\subsubsection{Matter-Only Perturbation (Fixed Metric)}
\label{sec:mfp-toml-matter-only}

\begin{lstlisting}[style=toml]
[linearization]
# perturbation_field omitted -> metric is NOT perturbed
# Internally uses: Perturbation[Bg[__], ___] := 0

[[linearization.matter_perturbations]]
field = "A"
perturbation_name = "a"
background = "Abar"
\end{lstlisting}

\subsubsection{Multiple Matter Fields}
\label{sec:mfp-toml-multiple}

\begin{lstlisting}[style=toml]
[linearization]
perturbation_field = "h"

[[linearization.matter_perturbations]]
field = "A"
perturbation_name = "a"
background = "Abar"

[[linearization.matter_perturbations]]
field = "phi"
perturbation_name = "dphi"
background = "phi0"
\end{lstlisting}

\subsubsection{Backward Compatibility}
\label{sec:mfp-toml-backward}

When \texttt{matter\_perturbations} is absent, behavior is identical to the current pipeline. No existing TOML configs are affected.

\subsection{Pipeline Architecture}
\label{sec:mfp-pipeline}

\subsubsection{Derived Fields and Perturbation Order}
\label{sec:mfp-derived-fields}

\textbf{Critical design detail}: The Lagrangian uses derived fields (e.g., $F$), but perturbation acts on fundamental fields ($A$).

\paragraph{Pipeline sequence.}
\label{par:mfp-pipeline-sequence}

\begin{enumerate}
  \item \textbf{User writes Lagrangian} with derived field $F$:
    \begin{equation*}
      \lag = -\tfrac{1}{4}\, F_{ab}\, \eta^{ac}\, \eta^{bd}\, F_{cd}
    \end{equation*}

  \item \textbf{MakeRule expansion} (existing \texttt{\_wls\_derived\_fields()} $\to$ \texttt{\_wls\_lagrangian()}):
    $F \to \mathrm{CD}_{-a}[A_{-b}] - \mathrm{CD}_{-b}[A_{-a}]$. The Lagrangian now has $\mathrm{CD}_{-a}[A_{-b}]$ terms.

  \item \textbf{xPert perturbation} of the expanded Lagrangian:
    \texttt{Perturbation[$\sqrtg\,\lag_\text{expanded}$, 2]} + \texttt{ExpandPerturbation[]}.
    xPert correctly produces:
    \begin{itemize}
      \item 0th order: $\mathrm{CD}_{-a}[A_{-b}]$ $\to$ background field strength $\bg{F}$
      \item 1st order: $\mathrm{CD}_{-a}[\texttt{aPert}[\texttt{LI}[1], -b]]$ $\to$ perturbation field strength $f$
    \end{itemize}

  \item \textbf{Background resolution} via xCoba during \texttt{ToBasis} + \texttt{TraceBasisDummy}:
    xCoba evaluates $\mathrm{CD}_{-a}[A_{-b}]$ $\to$ $\partial_a(\bg{A}_b)$ using known component values.
    For flat space: $\partial_z(\bg{A}_y) = \partial_z(-\Bzero z) = -\Bzero$.
\end{enumerate}

\textbf{Key point}: The existing \texttt{\_validate\_lagrangian} check rejects \texttt{CD[-a][G[]]} in the user's Lagrangian expression. This does NOT trigger because the user writes \texttt{F[-a,-b]} (not \texttt{CD[-a][A[-b]]}). The derivatives only appear after MakeRule expansion.

\subsubsection{Perturbation Order Truncation}
\label{sec:mfp-truncation}

TIDAL operates exclusively in the linearized regime. After \texttt{Perturbation[L, 2]} + \texttt{ExpandPerturbation}:

\begin{table}[htbp]
\centering
\caption{Perturbation order truncation rules.}
\label{tab:mfp-truncation}
\begin{tabular}{@{}llcl@{}}
\toprule
Term Type & Example & Keep? & Reason \\
\midrule
Products of 1st-order (same field) & $h[\texttt{LI}[1]] \times h[\texttt{LI}[1]]$ & \textbf{YES} & Quadratic action \\
Products of 1st-order (cross) & $h[\texttt{LI}[1]] \times a[\texttt{LI}[1]]$ & \textbf{YES} & Coupling terms \\
Pure 2nd-order & $h[\texttt{LI}[2]]$ or $a[\texttt{LI}[2]]$ & \textbf{NO} & Linearized: truncate at 1st order \\
\bottomrule
\end{tabular}
\end{table}

Dropping \texttt{LI[2]} is equivalent to setting field $=$ background $+ \varepsilon \cdot \text{perturbation}^{(1)}$ with no $\varepsilon^2$ terms. This is correct for all linearized theories.

\subsection{Wolfram Code Generation}
\label{sec:mfp-wolfram}

\subsubsection{New Function: \texttt{\_wls\_matter\_perturbation\_setup()}}
\label{sec:mfp-wls-setup}

Generates \texttt{DefTensorPerturbation} + \texttt{ComponentValue} for each matter perturbation.

\begin{lstlisting}[style=mathematica]
(* For each [[linearization.matter_perturbations]] entry: *)

(* 1. Define xPert perturbation for field A *)
DefTensorPerturbation[aPert[LI[order], -a], A[-a], M4];

(* 2. Set background values for A (0th order) *)
(* From matching [[background_fields]] entry *)
ComponentValue[A[{0, -Cart}], 0];
ComponentValue[A[{1, -Cart}], 0];
ComponentValue[A[{2, -Cart}], -B0 * z[]];
ComponentValue[A[{3, -Cart}], 0];
\end{lstlisting}

For tensor matter perturbations (rank-2):
\begin{lstlisting}[style=mathematica]
DefTensorPerturbation[TPert[LI[order], -a, -b], T[-a, -b], M4];
\end{lstlisting}

For scalar matter perturbations:
\begin{lstlisting}[style=mathematica]
DefTensorPerturbation[dphiPert[LI[order]], phi[], M4];
\end{lstlisting}

\subsubsection{Modified: \texttt{\_wls\_linearize\_from\_lagrangian()}}
\label{sec:mfp-wls-linearize}

After \texttt{ExpandPerturbation}, drop \texttt{LI[2]} and replace \texttt{LI[1]} for ALL perturbation fields:

\begin{lstlisting}[style=mathematica]
(* Existing: metric perturbation *)
SetupMetricPerturbation[g, hpert, epsilon];

(* NEW: each matter perturbation *)
DefTensorPerturbation[aPert[LI[order], -a], A[-a], M4];

(* 2nd-order perturbation of action density *)
l2Raw = Perturbation[Sqrt[-Detg[]] * Lagrangian, 2];
l2Raw = ExpandPerturbation[l2Raw];

(* Truncate ALL fields at 1st order *)
l2Raw = l2Raw /. hpert[LI[2], __] :> 0;
l2Raw = l2Raw /. aPert[LI[2], __] :> 0;

(* Replace xPert notation with declared field tensors *)
l2Raw = l2Raw /. hpert[LI[1], idx__] :> H[idx];
l2Raw = l2Raw /. aPert[LI[1], idx__] :> a[idx];

(* VarD for EACH dynamical field *)
eomH = VarD[H[-a,-b], CD][l2ForVarD];
eomA = VarD[a[-a], CD][l2ForVarD];
\end{lstlisting}

\subsubsection{Multi-Field VarD and Component Decomposition}
\label{sec:mfp-wls-decompose}

Each dynamical field gets its own \texttt{VarD} call. The results feed into \texttt{DecomposeToComponents} with all other dynamical fields listed as \texttt{additionalFields}:

\begin{lstlisting}[style=mathematica]
(* Decompose graviton equation -- a components are additional fields *)
DecomposeToComponents[eomH, H, {a}, ...]

(* Decompose photon equation -- H components are additional fields *)
DecomposeToComponents[eomA, a, {H}, ...]
\end{lstlisting}

The existing \texttt{additionalFields} mechanism in \texttt{DecomposeToComponents} already handles cross-field terms correctly (verified in \texttt{coupled\_scalars}, \texttt{scalar\_vector\_coupling} examples).

\subsection{Background Field Strength: How $\bg{F} = d\bg{A}$ Gets Evaluated}
\label{sec:mfp-background}

For the Gertsenshtein effect, the coupling coefficients contain the background field strength $\bg{F}$. Here is the evaluation chain:

\begin{enumerate}
  \item \textbf{TOML}: \texttt{[[background\_fields]]} with components \texttt{["0", "0", "-B0 * z[]", "0"]}
  \item \textbf{Wolfram}: \texttt{ComponentValue[A[\{2, -Cart\}], -B0 * z[]]}
  \item \textbf{After ExpandPerturbation}: 0th-order terms contain $\mathrm{CD}_{-a}[A_{-b}]$
  \item \textbf{ToBasis}: xCoba evaluates $\mathrm{CD}_{-a}[A_{-b}] \to \partial_a(\bg{A}_b)$ using known component values
  \item \textbf{Result}: For $\bg{A}_y = -\Bzero z$, xCoba computes $\partial_z(\bg{A}_y) = -\Bzero$
\end{enumerate}

For uniform $\Bzero$, the coupling coefficients are constants (proportional to $\Bzero$). For localized $\Bzero(z)$, they become position-dependent coefficients --- already supported by the solver's coordinate-dependent coefficient system (\texttt{IsCoordinateDependentCoefficient} $\to$ \texttt{\_mathematica\_to\_python()} $\to$ grid evaluation $\to$ L2 cache).

\subsection{Files Modified (Implemented)}
\label{sec:mfp-files}

\begin{table}[htbp]
\centering
\caption{Implementation status of multi-field perturbation pipeline.}
\label{tab:mfp-files}
\begin{tabular}{@{}lll@{}}
\toprule
File & Change & Status \\
\midrule
\texttt{tidal/cli/\_derive.py} & \texttt{\_validate\_matter\_perturbations()}, \texttt{\_validate\_single\_matter\_perturbation()} & Done \\
\texttt{tidal/cli/\_derive.py} & \texttt{\_wls\_matter\_perturbation\_setup()} --- DefTensorPerturbation + DefTensor & Done \\
\texttt{tidal/cli/\_derive.py} & \texttt{\_wls\_multi\_field\_eom()} --- multi-field VarD + DecomposeToComponents & Done \\
\texttt{tidal/cli/\_derive.py} & \texttt{\_wls\_matter\_pert\_truncation()} --- LI[2] drop + LI[1] replacement & Done \\
\texttt{tidal/cli/\_derive.py} & \texttt{\_xpert\_index\_pattern()}, \texttt{\_pert\_field\_dict()} --- helpers & Done \\
\texttt{tidal/cli/\_derive.py} & Extended \texttt{\_validate\_gauge()} for perturbation name targets & Done \\
\texttt{tidal/cli/\_derive.py} & Extended \texttt{\_wls\_linearize\_from\_lagrangian()} for multi-field path & Done \\
\texttt{tests/test\_cli.py} & \texttt{TestMatterPerturbations} class (6 tests) & Done \\
\texttt{examples/gertsenshtein/} & \texttt{theory.toml}, \texttt{run.sh}, \texttt{sweep\_coupling.sh} & Done \\
\texttt{tidal/wolfram/ComponentDecompose.wl} & No changes needed --- \texttt{additionalFields} already sufficient & Verified \\
\bottomrule
\end{tabular}
\end{table}

\subsection{Validation}
\label{sec:mfp-validation}

\begin{enumerate}
  \item \textbf{Regression}: All 1568 tests pass; 28 examples unaffected
  \item \textbf{TOML validation}: \texttt{matter\_perturbations} accepted/rejected correctly (5 error-case tests)
  \item \textbf{WLS generation}: \texttt{--dry-run} produces correct \texttt{DefTensorPerturbation} + multi-field \texttt{VarD} (verified)
  \item \textbf{End-to-end}: Einstein--Maxwell \texttt{--dry-run} generates correct WLS with coupling terms
  \item \textbf{Physics}: Conversion probability validated --- uniform $\Bzero$ (RMS~0.012, 40-point sweep) and localized Gaussian (0.04\% agreement with Boccaletti formula, 48-point sweep). See \texttt{docs/gertsenshtein.md}.
\end{enumerate}

\subsection{Design Principles}
\label{sec:mfp-design}

\begin{enumerate}
  \item \textbf{General-purpose}: Not Gertsenshtein-specific --- any theory with multiple perturbation fields
  \item \textbf{Lagrangian-first}: All equations derived from the action, never hardcoded
  \item \textbf{xPert-native}: Uses \texttt{DefTensorPerturbation}, not ad-hoc splitting
  \item \textbf{Backward-compatible}: No changes when \texttt{matter\_perturbations} is absent
  \item \textbf{Composable}: Works with existing gauge fixing, plane-wave reduction, background fields
\end{enumerate}
